\documentclass[10pt,a3paper]{article}
\usepackage[margin=10mm]{geometry}
\usepackage[T1]{fontenc}
\usepackage[utf8]{inputenc}
\usepackage[french]{babel}
\usepackage{lmodern}
\usepackage{microtype}
\makeatletter\def\input@path{{../../../Style/}}\makeatother
\NeedsTeXFormat{LaTeX2e}
\ProvidesPackage{TLPosterCarteMentale}[2026/08/03 Posters et cartes mentales SI]

\RequirePackage{tikz}
\RequirePackage[most]{tcolorbox}
\RequirePackage{xcolor}
\RequirePackage{amsmath,amssymb}
\RequirePackage{graphicx}
\usetikzlibrary{positioning,arrows.meta,calc,shapes.geometric,mindmap,backgrounds,fit}

\definecolor{TLPosterBlue}{RGB}{24,86,141}
\definecolor{TLPosterLight}{RGB}{234,243,250}
\definecolor{TLPosterAccent}{RGB}{232,126,34}
\definecolor{TLPosterGreen}{RGB}{52,152,102}
\definecolor{TLPosterRed}{RGB}{192,57,43}
\definecolor{TLPosterGray}{RGB}{80,90,100}

\newcommand{\TLPosterTitre}[2]{%
  \begin{tcolorbox}[enhanced,boxrule=0pt,arc=0pt,colback=TLPosterBlue,left=6mm,right=6mm,top=4mm,bottom=4mm]
    {\color{white}\bfseries\LARGE #1}\hfill{\color{white}\large #2}
  \end{tcolorbox}%
}

\newtcolorbox{TLPosterBloc}[2][]{
  enhanced,breakable,
  colback=TLPosterLight,
  colframe=TLPosterBlue,
  colbacktitle=TLPosterBlue,
  coltitle=white,
  fonttitle=\bfseries,
  title={#2},
  boxrule=0.8pt,
  arc=2mm,
  left=3mm,right=3mm,top=2mm,bottom=2mm,
  #1
}

\newtcolorbox{TLPosterFormule}[1][]{
  enhanced,
  colback=white,
  colframe=TLPosterAccent,
  boxrule=1pt,
  arc=2mm,
  fontupper=\bfseries,
  halign=center,
  #1
}

\newcommand{\TLPosterMethode}[1]{%
  \begin{tcolorbox}[enhanced,colback=TLPosterGreen!8,colframe=TLPosterGreen,title=Méthode,fonttitle=\bfseries]
  #1
  \end{tcolorbox}%
}

\newcommand{\TLPosterAttention}[1]{%
  \begin{tcolorbox}[enhanced,colback=TLPosterRed!6,colframe=TLPosterRed,title=Point de vigilance,fonttitle=\bfseries]
  #1
  \end{tcolorbox}%
}

\tikzset{
  TLmindroot/.style={concept,concept color=TLPosterBlue,text=white,font=\bfseries\large,minimum size=34mm},
  TLmindlevel1/.style={level 1 concept/.append style={font=\bfseries,minimum size=23mm,sibling angle=45}},
  TLmindlevel2/.style={level 2 concept/.append style={font=\small,minimum size=17mm}},
  TLarrow/.style={-{Stealth[length=3mm]},thick,draw=TLPosterBlue},
  TLbox/.style={rounded corners=2mm,draw=TLPosterBlue,fill=TLPosterLight,align=center,inner sep=3mm}
}

\newenvironment{TLCarteMentale}[1]{%
  \begin{tikzpicture}[mindmap,grow cyclic,concept color=TLPosterBlue,every node/.style={concept},TLmindlevel1,TLmindlevel2]
  \node[TLmindroot]{#1}
}{;
  \end{tikzpicture}
}

\newcommand{\TLBranche}[3][]{child[concept color=#2]{node[#1]{#3}}}

\endinput

\pagestyle{empty}
\renewcommand{\familydefault}{\sfdefault}
\begin{document}
\TLPosterCourseHeader{14}{Logique et systèmes à événements discrets}{Combinatoire, états, codeurs et communication}
\vspace{2mm}
\begin{minipage}[t]{0.61\linewidth}
\begin{TLPosterSavoirs}
\begin{itemize}
  \item distinguer combinatoire (sorties = fonction des entrées présentes) et séquentiel (sorties = entrées + état mémorisé)
  \item connaître NON, ET, OU, XOR, De Morgan, formes canoniques, Karnaugh et fonctions usuelles : multiplexeur, codeur, décodeur, comparateur, additionneur
  \item lire et construire un diagramme d'état : état, transition, événement, garde, effet, entry/do/exit, choix, fork/join et historique
  \item exploiter un codeur incrémental A/B/Z, le comptage en quadrature, un codeur absolu et le code Gray
  \item connaître débit, délai, BER, OSI, parité, CRC, trame, topologie et modes simplex/duplex, série/synchrone/asynchrone
\end{itemize}
\end{TLPosterSavoirs}
\end{minipage}\hfill
\begin{minipage}[t]{0.36\linewidth}
\TLPosterKey{Choisir le modèle}{Pas de mémoire $\Rightarrow$ table de vérité. Mémoire / enchaînement d'événements $\Rightarrow$ graphe d'états.}
\end{minipage}
\vfill
\begin{minipage}[t]{0.49\linewidth}
\begin{TLPosterBloc}[Méthode combinatoire]{TLPosterBlue}
1) Nommer entrées/sorties et leur état actif. 2) Construire la table de vérité, avec états indifférents si besoin. 3) Écrire une forme canonique à partir des lignes où la sortie vaut 1. 4) Simplifier par Boole ou Karnaugh : groupes de $1,2,4,8\ldots$, les plus grands possibles, bords opposés adjacents. 5) Réaliser le logigramme et vérifier toutes les combinaisons utiles.
\end{TLPosterBloc}
\begin{TLPosterBloc}[Méthode diagramme d'état]{TLPosterPurple}
1) Recenser les situations stables. 2) Associer activités \texttt{entry/do/exit}. 3) Tracer les transitions sous la forme événement [garde] / effet. 4) Ajouter état initial, choix, parallélisme ou temporisations \texttt{after(T)} si nécessaire. 5) Vérifier : tous les cas couverts, gardes non contradictoires, priorités de sécurité explicites et absence d'état inaccessible.
\end{TLPosterBloc}
\end{minipage}\hfill
\begin{minipage}[t]{0.49\linewidth}
\begin{TLPosterBloc}[Méthode codeur]{TLPosterGreen}
\textbf{Incrémental :} $N$ périodes/tour ; sur 4 fronts, $N_c=4N$ et $\Delta\theta=2\pi/N_c$. A en avance sur B $\Rightarrow$ un sens ; B en avance sur A $\Rightarrow$ sens inverse. Position par comptage ; vitesse par nombre d'impulsions sur une fenêtre à haute vitesse, ou période entre fronts à basse vitesse.\\
\textbf{Absolu :} $n$ pistes $\Rightarrow 2^n$ positions ; utiliser Gray pour éviter plusieurs changements de bits simultanés.
\end{TLPosterBloc}
\begin{TLPosterBloc}[Méthode trame / réseau]{TLPosterTeal}
\textbf{Parité :} compter les bits à 1 et vérifier la règle paire/impaire.\\
\textbf{CRC :} données $\rightarrow$ ajouter $r$ zéros $\rightarrow$ division modulo 2 par $G(X)$ $\rightarrow$ remplacer par le reste $\rightarrow$ vérifier reste nul à la réception.\\
\textbf{Réseau :} identifier débit utile, délai/déterminisme, support, topologie, mode d'échange et mécanisme de détection/retransmission avant de juger une solution.
\end{TLPosterBloc}
\end{minipage}
\vfill
\begin{TLPosterMethode}
\begin{enumerate}
  \item \textbf{Si fonction logique sans mémoire :} cahier des charges $\rightarrow$ table de vérité $\rightarrow$ équation $\rightarrow$ simplification $\rightarrow$ implantation $\rightarrow$ test.
  \item \textbf{Si comportement séquentiel :} états stables $\rightarrow$ transitions $\rightarrow$ événements/gardes/effets $\rightarrow$ temporisations/parallélisme $\rightarrow$ vérification de complétude.
  \item \textbf{Si codeur :} identifier incrémental ou absolu, calculer la résolution, puis choisir la méthode de position/vitesse adaptée.
  \item \textbf{Si communication :} lire la trame dans l'ordre, séparer données et contrôle, vérifier parité/CRC puis raisonner sur débit, délai, fiabilité et déterminisme.
  \item \textbf{Toujours :} traiter explicitement les états impossibles, défauts, priorités et transitions invalides.
\end{enumerate}
\end{TLPosterMethode}
\vfill
\TLPosterRetenir{commencer par décider si le problème est combinatoire, séquentiel, de mesure numérique ou de communication : chaque cas a sa méthode propre.}
\end{document}
